\section{The stability of the non-rotating 2D SWE  on a B-grid}

\textbf{(a)} The non-rotating 2-D SWE equations are given as 
\begin{subequations}
\begin{empheq}[left=\empheqlbrace]{align}
\frac{\partial u}{\partial t}& = -g \frac{\partial h}{\partial x} \\
\frac{\partial v}{\partial t}& = -g \frac{\partial h}{\partial y} \\
\frac{\partial h}{\partial t}& + H \left( \frac{\partial u}{\partial x} + \frac{\partial v}{\partial y} \right) = 0.
\label{eq:SWE}
\end{empheq}
\end{subequations}
Discretising the equations on an Arakawa B-grid by applying the leapfrog scheme one obtains
\begin{subequations}
\begin{empheq}[left=\empheqlbrace]{align}
\frac{u_{i,j}^{n+1} - u_{i,j}^{n-1}}{2\Delta t} &= -g \frac{h_{i+1,j+1}^{n} + h_{i+1,j}^{n}-h_{i,j+1}^{n}-h_{i,j}^{n}}{2\Delta x} \\
\frac{v_{i,j}^{n+1} - v_{i,j}^{n-1}}{2\Delta t} &= -g \frac{h_{i,j+1}^{n} + h_{i+1,j+1}^{n}-h_{i+1,j}^{n}-h_{i,j}^{n}}{2\Delta y}
\\
\frac{h_{i,j}^{n+1} - h_{i,j}^{n-1}}{2\Delta t}& = - H \left( \frac{u_{i,j}^{n} + u_{i,j-1}^{n}-u_{i-1,j}^{n}-u_{i-1,j-1}^{n}}{2\Delta x} + \frac{v_{i,j}^{n} + v_{i-1,j}^{n}-v_{i,j-1}^{n}-v_{i-1,j-1}^{n}}{2\Delta y} \right).
\end{empheq}
\label{eq:discretised_SWE}
\end{subequations}

\textbf{(b)} To solve for the stability criterion we apply a von Neumann stability analysis. We assume that the solution can be written as a Fourier mode of the form
\begin{equation}\begin{pmatrix}
u_{i,j}^{n} \\
v_{i,j}^{n} \\
h_{i,j}^{n}
\end{pmatrix} = \begin{pmatrix}u_0 \\
v_0 \\
h_0
\end{pmatrix} e^{\mathfrak{I}(k i \Delta x + l j \Delta y - \omega n \Delta t)},
\end{equation}
where $k$ and $l$ are the wavenumbers in the $x$ and $y$ direction, respectively, $\omega$ is the frequency and $\mathfrak{I}$ is the imaginary unit. 

One can simplify \autoref{eq:discretised_SWE} as
\begin{subequations}
\begin{align*}
u_0 &= -\frac{1}{\lambda - \lambda^{-1}} \frac{g h_0 \Delta t}{\Delta x} \left(e^{\mathfrak{I}k \Delta x}e^{\mathfrak{I}l \Delta y}+ e^{\mathfrak{I}k\Delta x} - e^{\mathfrak{I}l \Delta y}-1\right) \\
v_0 &= -\frac{1}{\lambda - \lambda^{-1}} \frac{gh_0\Delta t}{\Delta y} \left(e^{\mathfrak{I}k \Delta x}e^{\mathfrak{I}l \Delta y}+ e^{\mathfrak{I}l\Delta y} - e^{\mathfrak{I}k \Delta x}-1\right) \\
\begin{split}
    \left(\lambda - \lambda^{-1}\right) h_0 = &- \frac{Hu_0\Delta t}{\Delta x}  \left(1+e^{-\mathfrak{I}l \Delta y} - e^{-\mathfrak{I}k \Delta x}- e^{-\mathfrak{I}k \Delta x}e^{-\mathfrak{I}l \Delta y}\right)  \\ 
    &- \frac{Hv_0\Delta t}{\Delta y}\left(1+e^{-\mathfrak{I}k \Delta x} - e^{-\mathfrak{I}l \Delta y}- e^{-\mathfrak{I}k \Delta x}e^{-\mathfrak{I}l \Delta y}\right).
\end{split}
\end{align*}
\end{subequations}
One can simplify this by assuming $\Delta x = \Delta y$ and defining $\mu = \sqrt{gH}\,\Delta t/\Delta x$. Then inserting the expressions for $u_0$ and $v_0$ into the last equation gives
\begin{equation*}
    \begin{split}
    \left(\lambda - \lambda^{-1}\right)^2 &= \mu^2 \left(e^{\mathfrak{I}(k+l) \Delta x}+ e^{\mathfrak{I}k\Delta x} - e^{\mathfrak{I}l \Delta x}-1\right)\left(1+e^{-\mathfrak{I}l \Delta x} - e^{-\mathfrak{I}k \Delta x}- e^{-\mathfrak{I}(k+l) \Delta x}\right) \\ 
    &
    + \mu^2 \left(e^{\mathfrak{I}(k+l) \Delta x}+ e^{\mathfrak{I}l\Delta x} - e^{\mathfrak{I}k \Delta x}-1\right)\left(1+e^{-\mathfrak{I}k \Delta x} - e^{-\mathfrak{I}l \Delta x}- e^{-\mathfrak{I}(k+l) \Delta x}\right).
    \end{split}
\end{equation*}
Here one can make the neat observation by using the complex conjugates
\begin{equation*}
\begin{split}
(\lambda - \lambda^{-1})^2
&= \mu^2 \,(e^{\mathfrak{I}k\Delta x}-1)(e^{\mathfrak{I}l\Delta x}+1)
\,(1 - e^{-\mathfrak{I}k\Delta x})(1 + e^{-\mathfrak{I}l\Delta x}) \\
&\quad
+ \mu^2 \,(e^{\mathfrak{I}l\Delta x}-1)(e^{\mathfrak{I}k\Delta x}+1)
\,(1 - e^{-\mathfrak{I}l\Delta x})(1 + e^{-\mathfrak{I}k\Delta x}).
\end{split}
\end{equation*}
Now using the mathematical fact that 
\begin{equation}
    \left(e^{\mathfrak{I}\theta} - 1\right)\left(e^{-\mathfrak{I}\theta}-1\right) = 4\sin^2\left(\frac{\theta}{2}\right) \quad \text{and} \quad \left(e^{\mathfrak{I}\theta} + 1\right)\left(e^{-\mathfrak{I}\theta}+1\right) = 4\cos^2\left(\frac{\theta}{2}\right),
\end{equation}
one can simplify the above expression to
\begin{equation*}
(\lambda - \lambda^{-1})^2
= -16\mu^2 \left[
\sin^2\left(\frac{k\Delta x}{2}\right)\cos^2\left(\frac{l\Delta x}{2}\right)
+
\cos^2\left(\frac{k\Delta x}{2}\right)\sin^2\left(\frac{l\Delta x}{2}\right)
\right].
\end{equation*}
We now simply define the positive trigonometric expression as $\mathcal{A}^2$. 
Taking the square root of the above expression gives
\begin{equation*}
\lambda - \lambda^{-1} = \pm 4\mathcal{A}\mu \mathfrak{I}.
\end{equation*}
Which gives the following expression for $\lambda$:
\begin{equation*}
\lambda = \pm 2\mathcal{A}\mu \mathfrak{I} \pm \sqrt{1 - 4\mathcal{A}^2\mu^2}.
\end{equation*}
To evaluate the stability of the scheme we need to evaluate the magnitude of $\lambda$. We do this by considering if the square root is real or imaginary. If the square root is imaginary, then the magnitude of $\lambda$ is given by
\begin{equation*}
|\lambda|^2 = 4\mathcal{A}^2\mu^2 + \text{root}^2.
\end{equation*}
This is greater than 1 for any value of $\mu$ which implies that the root must be real for the scheme to be stable. Thus if the root is real, then the magnitude of $\lambda$ is given by
\begin{equation*}
|\lambda|^2 = 4\mathcal{A}^2\mu^2 + \left(1 - 4\mathcal{A}^2\mu^2\right) = 1.
\end{equation*}
Thus the scheme is stable if the root is real, which implies that
\begin{equation*}
    1 - 4\mathcal{A}^2\mu^2 \geq 0 \quad \Rightarrow \quad \mu \leq \frac{1}{2\mathcal{A}}.
\end{equation*}
And since $\mathcal{A}^2$ is positive and has a maximum value of 1, we can thus conclude that the stability criterion. 

\begin{answer}
    The stability criterion for the non-rotating 2D SWE on a B-grid is given by $\mu \leq \frac{1}{2\mathcal{A}}$, where $\mu$ is the Courant number and $\mathcal{A}$ is a positive trigonometric expression defined as $\mathcal{A}^2 = \sin^2\left(\frac{k\Delta x}{2}\right)\cos^2\left(\frac{l\Delta x}{2}\right) + \cos^2\left(\frac{k\Delta x}{2}\right)\sin^2\left(\frac{l\Delta x}{2}\right)$. 
    \\ $ $ \\
    This implies that the scheme is stable if the Courant number $\mu \leq \frac{1}{2}$.
\end{answer}